\section{The stationary minimax exponent}
\label{sec:sharp-stationary}
The stationary upper and lower powers can be matched. The geometric
point is a cancellation between the two curved sides of a collision
strip. Only boundary arcs for which one translated endpoint is in the
launch disk and the other is outside can contribute to the derivative
of a collision probability. Each such arc, and each swept facet, forces
a cubic lower bound for both the collision area and its complement.
This gives a Hellinger contraction of order $\epsilon^{3/2}$ uniformly
in the command length, including lengths tending to zero. The conditional
information bound and the physical packing then give the stationary
lower power. The launch disk remains fixed and its law may be disclosed
in full; the result does not concern a noise radius tending to zero.

\subsection{A uniform Hellinger bound for short commands}
The cancellation between the physical and swept boundaries remains
effective when the command length tends to zero. We establish the
geometric estimate needed to quantify this cancellation. Throughout this
subsection, $n(\theta)=(\cos\theta,\sin\theta)$ and
$\tau(\theta)=(-\sin\theta,\cos\theta)$.

\begin{lemma}[Effective collision boundary]
\label{lem:effective-collision-boundary}
There is a constant $R_0>0$ with the following property. Fix
$0<R_-\le R_+\le R_0$. For some $\zeta_*>0$, let $C$ be a smooth convex
body whose support function satisfies
\[
                         \|p_C-1\|_{C^2}\le\zeta_*.
\]
For $R\in[R_-,R_+]$, a disk $D=x+R\overline B$, a unit vector $e$,
and $t\ge0$, put
\[
 S=C+[-te,0],\qquad E=S\setminus C,\qquad
 A=\area(D\cap E),\qquad A_c=\area(D)-A.
\]
Write $b=p_Cn+p_C'\tau$, $r=p_C+p_C''$, and
$H_e=\{\theta:n(\theta)\cdot e<0\}$. Let $F_+$ and $F_-$ be the
two facets of $S$ with normals $e^\perp$ and $-e^\perp$, respectively.
Then
\begin{align}
 L_D(C;t,e)
 &: =\int_{H_e}r(\theta)
       \left|\one_D(b(\theta)-te)-\one_D(b(\theta))\right|\dd\theta
       +\mathcal H^1(F_+\cap D)+\mathcal H^1(F_-\cap D)
       \label{eq:effective-boundary-length}\\
 &\le K\min\{A,A_c\}^{1/3}.\label{eq:effective-boundary-area}
\end{align}
Here $K$ is independent of $x,e,t,R$, and $C$ in the stated class.
\end{lemma}
\begin{proof}
Choose $R_0$ sufficiently small and then $\zeta_*$ sufficiently small
relative to $R_-$. All curvature radii lie in $[1/2,2]$. Rotate the
coordinates so that $e=(1,0)$, and parametrize the entering boundary by
arclength $s$. The map
\[
               (s,q)\longmapsto b(s)-qe,\qquad 0<q<t,
\]
parametrizes $E$ injectively off its boundary. Its area Jacobian is
$a(s)=|n(s)\cdot e|$. This follows also by slicing horizontally:
if the left endpoint of $C$ at height $y$ is $l(y)$, the corresponding
collision interval is $[l(y)-t,l(y))$.

We first bound the number of boundary pieces in
\eqref{eq:effective-boundary-length}. The intersection of $\partial C$
with any disk of the stated radius is empty or a single short arc.
Indeed, $b(\theta)=n(\theta)+O(\zeta_*)$. If $b(\theta)\in D$, its
angle lies in an interval of length $O(R+\zeta_*)$ about the direction
of $x$. The whole candidate arc is at distance $O(R+\zeta_*)$ from
$x$. On that arc, with derivatives taken in arclength,
\[
 \frac{\dd^2}{\dd s^2}|b(s)-x|^2
       =2+2(b(s)-x)\cdot b''(s)\ge1.
\]
Its disk sublevel set is therefore an interval. Intersecting this short
arc with the entering hemisphere preserves this property. Applying the
same argument with center $x+te$ shows that the two disk-membership sets
for $b$ and $b-te$ are intervals. Their symmetric difference has at most
two interval components. Each facet contributes at most one chord.
Thus $L_D$ is the sum of at most four lengths.

Consider one active arc $I$, of arclength $\ell$. Its endpoint lying
in $D$ is denoted by $z(s)$. Either $z=b$ and the direction toward the
other endpoint is $-e$, or $z=b-te$ and that direction is $e$.
Write it as $\sigma e$, where $\sigma\in\{-1,1\}$; then
\[
                z(s)\in D,\qquad z(s)+\sigma te\notin D
                         \quad(s\in I),
\]
apart from endpoints of measure zero. Let $J$ be the middle half of
$I$. Its points are at arclength distance at least $\ell/4$ from
either endpoint of the entering hemisphere. The curvature bounds give
\begin{equation}\label{eq:active-entering-angle}
                         a(s)\ge c\ell\qquad(s\in J).
\end{equation}
For such $s$, set
\[
 v=z(s)-x,\qquad B=|v\cdot\tau(s)|,\qquad G=R^2-|v|^2.
\]
The squared-distance function $f(s)=|z(s)-x|^2$ satisfies $f''\ge1$
on $I$. Its values at $s\pm\ell/8$ do not exceed $R^2$, whereas
$f'(s)=2v\cdot\tau(s)$. The two lower Taylor estimates imply
\begin{equation}\label{eq:active-disk-clearance}
                              G\ge c\ell(B+\ell).
\end{equation}
Let $q(s)$ be the distance to the disk boundary along the ray
$z(s)+q\sigma e$. The other endpoint is outside $D$, so $q(s)\le t$.
The disk equation gives
\[
 q(s)=\frac{G}{\sqrt{(v\cdot e)^2+G}+\sigma v\cdot e}
       \ge\frac{G}{2|v\cdot e|+\sqrt G}.
\]
Since $|v\cdot e|\le Ra(s)+B$, the monotonicity of
$G/(c+\sqrt G)$ in $G$, together with
\eqref{eq:active-disk-clearance}, yields
\[
 q(s)\ge c\frac{\ell(B+\ell)}{a(s)+B+\ell},\qquad
                         a(s)q(s)\ge c\ell^2.
\]
For the last inequality use \eqref{eq:active-entering-angle}, considering
$B\le a(s)$ and $B\ge a(s)$. The ray up to $q(s)$ lies in the
collision strip in either of the two activity cases. Its injective
parametrization therefore gives
\begin{equation}\label{eq:active-arc-volume}
                A\ge\int_Ja(s)q(s)\dd s\ge c\ell^3.
\end{equation}
In particular, no division by $t$ is involved.

Next let $I=F_+\cap D$ have length $\ell$, and let $J$ again be its
middle half. The squared distance along the facet has second derivative
$2$. The preceding Taylor argument shows that the Euclidean disk
clearance on $J$ is at least $c\ell^2$. Moreover, $J$ is at distance
at least $\ell/4$ from both endpoints of the whole facet. At inward
depth $w$ below its supporting line, the horizontal displacement of
the entering boundary from its tangency point is at most $C\sqrt w$.
Indeed, the curvature-radius bounds make the vertical drop comparable
to the square of angular distance and the horizontal displacement
comparable to angular distance. For $0\le w\le c\ell^2$, choose $c$
so that this displacement is at most $\ell/8$. The horizontal collision
interval then covers $J$. A rectangle of length $\ell/2$ and inward
depth $c\ell^2$ lies in $D\cap E$, proving
\eqref{eq:active-arc-volume} for the facet chord. The bottom facet has
the identical argument with the inward direction reversed.

We also need the complementary estimate. On the middle half of any
active arc, \eqref{eq:active-disk-clearance} gives Euclidean clearance
\[
                     R-|v|=\frac{G}{R+|v|}\ge c\ell^2.
\]
If $z=b$, the normal strip of inward depth $c\ell^2$ lies in $C$.
If $z=b-te$, the normal strip of outward depth $c\ell^2$ lies outside
$S$, since $b-te$ is its support point with that normal. Both strips
lie in $D\setminus E$. Their normal-coordinate Jacobians are bounded
below uniformly. To justify inward injectivity explicitly, choose a
fixed $q_*>0$ smaller than every curvature radius. The function
$p_C-q_*$ is a support function, so $C=C_0+q_*\overline B$; inward
offsets of depth less than $q_*$ are disjoint parallel boundaries.
Outward injectivity follows from uniqueness of projection onto the
convex set $S$. These strips prove $A_c\ge c\ell^3$.
For a facet chord, its outward rectangle of the same depth lies outside
$S$ and inside $D$, proving the same estimate. Each of the at most four
effective lengths is consequently bounded by
$K\min\{A,A_c\}^{1/3}$. Their sum proves
\eqref{eq:effective-boundary-area}. When $t=0$, both sides are zero.
\end{proof}

\begin{lemma}[Uniform short-command Hellinger contraction]
\label{lem:all-short-hellinger}
Let $C_0,C_1$ belong to the class in
Lemma~\ref{lem:effective-collision-boundary}, and put
$\epsilon=\|p_{C_1}-p_{C_0}\|_\infty$. For the same $D,t,e$, let
\[
 P_i=\frac{\area\bigl(D\cap((C_i+[-te,0])\setminus C_i)\bigr)}
                 {\pi R^2},\qquad i\in\{0,1\}.
\]
Then, uniformly in the nominal center, direction, command length
$t\ge0$, and radius $R\in[R_-,R_+]$,
\begin{equation}\label{eq:all-short-hellinger}
 H^2(\operatorname{Ber}(P_0),\operatorname{Ber}(P_1))
 :=(\sqrt{P_0}-\sqrt{P_1})^2
    +(\sqrt{1-P_0}-\sqrt{1-P_1})^2
                         \le K\epsilon^{3/2}.
\end{equation}
\end{lemma}
\begin{proof}
Interpolate the support functions by
$p_\lambda=(1-\lambda)p_{C_0}+\lambda p_{C_1}$, $0\le\lambda\le1$,
and put $u=p_{C_1}-p_{C_0}$, $r_\lambda=p_\lambda+p_\lambda''$,
$b_\lambda=p_\lambda n+p_\lambda'\tau$. These bodies remain in the
same fixed neighborhood. Write $A(\lambda)$ for their collision areas
in $D$. The normal velocity of both the physical boundary and its
translate is $u(\theta)$. The facets have the corresponding constant
normal velocities. Differentiating the swept area minus the physical
area gives, for almost every $\lambda$,
\begin{align}
 A'(\lambda)
 &=\int_{H_e}u(\theta)r_\lambda(\theta)
     \bigl[\one_D(b_\lambda(\theta)-te)
                  -\one_D(b_\lambda(\theta))\bigr]\dd\theta
          \label{eq:collision-shape-derivative}\\
 &\quad+u(e^\perp)\mathcal H^1(F_{+,\lambda}\cap D)
       +u(-e^\perp)\mathcal H^1(F_{-,\lambda}\cap D).\nonumber
\end{align}
Here $u(\pm e^\perp)$ denotes evaluation at the corresponding normal.
For completeness, the parallel-area bound makes $A$ absolutely
continuous in $\lambda$. The moving curved boundaries have curvature
strictly smaller than that of the disk, and their disk intersections
are the finite arcs described above. Their intersections with
$\partial D$ therefore have arclength zero, as do those of the facets.
The moving-boundary strip formula, followed by dominated convergence,
justifies \eqref{eq:collision-shape-derivative}. The nonentering
curved boundaries cancel because their normal velocities and arclength
densities agree. The entering terms likewise cancel whenever both
corresponding points belong to $D$.

Since $\|u\|_\infty=\epsilon$, the preceding lemma implies, for
$P(\lambda)=A(\lambda)/(\pi R^2)$,
\begin{equation}\label{eq:collision-probability-differential}
 |P'(\lambda)|\le K\epsilon
                   \min\{P(\lambda),1-P(\lambda)\}^{1/3}.
\end{equation}
This inequality includes probabilities arbitrarily close to either
endpoint. Apply the ordinary absolutely continuous chain rule to
$(P+\delta)^{2/3}$ and $(1-P+\delta)^{2/3}$, where $\delta>0$.
Equation~\eqref{eq:collision-probability-differential} bounds each
derivative by $K\epsilon$. Integrate from $0$ to $1$ and let
$\delta\downarrow0$. Thus
\[
 |P_1^{2/3}-P_0^{2/3}|\le K\epsilon,\qquad
 |(1-P_1)^{2/3}-(1-P_0)^{2/3}|\le K\epsilon.
\]
For nonnegative $a,b$, $|a^{3/4}-b^{3/4}|\le|a-b|^{3/4}$.
Applying this inequality to the two preceding estimates, and then
squaring, proves \eqref{eq:all-short-hellinger}. At $t=0$ the
probabilities are both zero.
\end{proof}

The constants are uniform on the stated positive interval of launch
radii. Parameter-independent mixtures of directions preserve
\eqref{eq:all-short-hellinger} by joint convexity of squared Hellinger
distance. A reverse command only translates the nominal disk and
reverses the direction, both already uniform in the lemma.

\subsection{Binary information and physical reconstruction}
\begin{lemma}[Hellinger diameter bounds binary information]
\label{lem:hellinger-capacity}
Let $V$ have any distribution on a finite set and, conditional on
$V=v$, let $Y$ be Bernoulli with mean $p_v$. Put
$a=\min_vp_v$ and $b=\max_vp_v$. Then, in bits,
\begin{equation}\label{eq:hellinger-capacity}
 I(V;Y)\le\frac1{\ln2}
 H^2(\operatorname{Ber}(a),\operatorname{Ber}(b)).
\end{equation}
No positive lower bound on $p_v$ or $1-p_v$ is required. The same
assertion holds under every posterior distribution.
\end{lemma}
\begin{proof}
Put $q=(a+b)/2$. If $q$ is zero or one, the observation is constant.
Otherwise comparison with $Q=\operatorname{Ber}(q)$ gives
\[
 I(V;Y)\le\mathbb E D_2(\operatorname{Ber}(p_V)\Vert Q)
 \le\frac{(b-a)^2}{4\ln2\,q(1-q)}.
\]
Here $D_2$ denotes relative entropy with logarithm to base two.
The first inequality follows by subtracting the nonnegative divergence
of the mixture law from $Q$; the second follows from
$\ln u\le u-1$, which bounds relative entropy by chi-square
divergence. On the other hand,
\begin{align*}
 H^2(\operatorname{Ber}(a),\operatorname{Ber}(b))
 &=(b-a)^2\left\{
 \frac1{(\sqrt a+\sqrt b)^2}
 +\frac1{(\sqrt{1-a}+\sqrt{1-b})^2}\right\}\\
 &\ge\frac{(b-a)^2}{4q(1-q)}.
\end{align*}
Here $(\sqrt a+\sqrt b)^2\le2(a+b)=4q$, and the
complementary inequality is identical. Endpoint cases follow by
continuity. Nothing in the proof restricts the distribution of $V$.
\end{proof}

\begin{theorem}[The stationary lower power for arbitrary short commands]
\label{thm:sharp-stationary-lower}
For every $s=6+\beta$, $0<\beta\le1$, there are fixed physical
priors, a nondegenerate subclass $\mathcal P\subset\mathcal B_\eta$,
a fixed $t_{\max}>0$, and a fixed known uniform-disk launch law
with the following property. Every admissible adaptive controller using
fresh draws of that law, exact nominal centers, and collision bits from
arbitrary directions and lengths $|a|\le t_{\max}$, and achieving
centered component $C^2$ error at most $\nu$ with probability at least
$1-\delta$ uniformly on $\mathcal P$, satisfies
\begin{equation}\label{eq:sharp-stationary-lower}
 \sup_{\mathcal O\in\mathcal P}\mathbb E_{\mathcal O}\mathsf T
                  \ge c\nu^{-(3s/2+1)/(s-2)}
\end{equation}
for fixed $0<\delta<1/4$ and all sufficiently small $\nu$.

The same assertion holds if the controller chooses the known disk
radius in a fixed interval $0<R_-\le R\le R_+$. The interval and
the physical priors are chosen once, with
$t_{\max}+2R_+<d_0$ and with $R_+$ small enough for
Lemma~\ref{lem:all-short-hellinger}. The lattice, primitive orbit
count, species, component correspondences, and parameter-independent
laboratory frame may be disclosed. The actual launch points remain
unobserved. Continuous adaptive directions and arbitrarily small
positive command lengths are included.
\end{theorem}
\begin{proof}
Use the fixed lattice, fixed second obstacle, and support-bump family
of Lemmas~\ref{lem:packing} and \ref{lem:centered-packing}.
Choose the fixed near-circle neighborhood, launch-radius interval, and
strict separation slack first. For all small $h$ the packing lies in
that neighborhood and contains $M=2^m$ tables such that
\begin{align*}
 m&\ge c_0h^{-1},\qquad
 \|p^0_\omega-p^0_{\omega'}\|_{C^2}
                  \ge c_1h^{s-2}\quad(\omega\ne\omega'),\\
 \max_{\omega,\omega'}\|p_\omega-p_{\omega'}\|_\infty
                  &\le C_0h^s.
\end{align*}
The $C^2$ distance from the unit disk is $O(h^{s-2})$.
The laboratory frame and all disclosed data are the same for every
packing parameter; no unknown Steiner point is disclosed.

We specify the component reduction uniformly over the packing.
For a nominal disk $D_R(x)$ and a displacement $a$, every possible
flight lies in the capsule $D_R(x)+[0,a]$, whose diameter is at
most $2R_++t_{\max}$. Relative to one packing member, every other
member's indexed component is within Hausdorff distance $C_0h^s$.
Enlarge the capsule by this amount. For small $h$ its diameter is
still strictly less than $d_0$, so it meets at most one component
of the reference table. Consequently the original capsule can meet
only one common indexed component family across all parameters.
It may miss that family for some parameters. If that component is
fixed, the response law is parameter-independent. If it is a variable
component, Lemma~\ref{lem:all-short-hellinger}, after the common
lattice translation, applies to every pair. Thus every available
Bernoulli response family has Hellinger diameter at most
\begin{equation}\label{eq:sharp-per-attempt-hellinger}
                             C h^{3s/2}.
\end{equation}
All constants are uniform in the nominal center, direction, length,
and chosen launch radius. A reciprocal reverse command is the same
kind of command after translating its nominal center and reversing
its displacement. A parameter-independent mixture of directions
preserves the bound by joint convexity of squared Hellinger distance;
alternatively one may reveal the direction before discarding it.

Condition on the independent controller seed. At every nonterminal
history the posterior is supported on the same finite packing, and
the selected command is fixed by that history and seed.
Lemma~\ref{lem:hellinger-capacity} and
\eqref{eq:sharp-per-attempt-hellinger} bound its conditional
information by $C h^{3s/2}$. Pad the binary record with a fixed
symbol after stopping. The finite chain rule gives
\[
 I(V;Z_{\le n},U_0)
     \le C h^{3s/2}\sum_{k=1}^n\Prob\{\mathsf T\ge k\}
     = C h^{3s/2}\mathbb E\min(n,\mathsf T).
\]
The conditional-entropy convergence and monotone-convergence argument
of Proposition~\ref{prop:stopped-contraction} gives the same
inequality for the complete stopped record, with
$\mathbb E\mathsf T$ on the right. Infinite expected cost already
satisfies the theorem.

Take $h$ to be a sufficiently large fixed multiple of
$\nu^{1/(s-2)}$. The centered-shape accuracy sets are disjoint.
Fano's inequality in its stated stopped-transcript form yields
\[
 C h^{3s/2}\frac1M\sum_\omega\mathbb E_\omega\mathsf T
    \ge (1-\delta)m-h_2(\delta)\ge c h^{-1}.
\]
The supremum dominates the average, proving
\eqref{eq:sharp-stationary-lower}. The fixed-radius assertion is
the special case in which the controller uses one radius throughout.
\end{proof}

\subsection{Shrinking boundary layers and the attempted-bit logarithm}
The fixed-layer bisection of Theorem~\ref{thm:rare-stationary} uses the
final probability accuracy at every level. It remains a valid controller.
The following additional result uses a coarser layer while the radial
interval is wide. A safeguarded interval update permits the layer to
shrink without losing the true boundary from the interval.

\begin{proposition}[A shrinking-layer rare-collision controller]
\label{prop:shrinking-rare}
Under the hypotheses of Theorem~\ref{thm:rare-stationary}, its geometric
and primitive-data conclusions can be achieved with
\begin{align}
 J_\nu&\le C\nu^{-1/(s-2)}\log(C/\nu),
                                  \label{eq:shrinking-centers}\\
 N_\nu&\le C\nu^{-((\gamma+3/2)s+1)/(s-2)}
                                      \log(C/(\nu\delta)).
                                  \label{eq:shrinking-attempts}
\end{align}
The assertion includes both the calibrated one-scale experiment and
the two known scales with unknown footprint, with their respective
fixed compass-separation hypotheses. Nominal centers require only the
same finest precision $O(\nu^{s/(s-2)})$. The launch spread, direction
law and fixed aperture are unchanged.
\end{proposition}
\begin{proof}
Reserve failure probability $\delta/2$ for the fixed-accuracy acquisition
of Lemma~\ref{lem:rare-coarse-normals}. We first describe one radial
search on this successful coarse event. Translate its interior center
to the origin, and write
\[
 \overline B_r\subset P\subset\overline B_{R_1},
 \qquad R=R_P(u)
\]
for its known uniform radial bounds and unknown boundary radius in the
chosen unit direction $u$. Fix a dyadic constant $A\ge4R_1/r$.
For all sufficiently small $e>0$, the inball inclusions
\begin{equation}\label{eq:shrinking-radial-conversion}
 (1-\rho/r)P\subset P_{-\rho},\qquad
 P+\rho\overline B\subset(1+\rho/r)P
                      \quad(0<\rho<r)
\end{equation}
show the following. If a nominal point within distance $e/2$ of $au$
is used in Proposition~\ref{prop:rare-query} with physical layer $e$,
its label is correct whenever $|a-R|>Ae$, except on the conditional
failure event of that proposition. Indeed $a<R-Ae$ places $au$ in
$P_{-2e}$ by the first inclusion with $\rho=2e$, whereas
$a>R+Ae$ places it outside $P+2e\overline B$ by the second.
The rounding reserve preserves depth at least $e$ in the first case
and an exterior point in the second. Thus, on a successful test,
\begin{equation}\label{eq:shrinking-label-invariant}
 \text{label }1\ \Longrightarrow\ a\le R+Ae,
 \qquad
 \text{label }0\ \Longrightarrow\ a\ge R-Ae.
\end{equation}
The fixed tubular neighborhood of the original coarse bracket supplies
the same valid candidate normal for these rounded points.

Take an initial certified interval $[l_0,u_0]$ containing $R$, with
fixed width $w_0>0$. Coarse accuracy and numerical reserves may be
chosen so that this width is a common dyadic prior constant for all
radial searches, and so that $w_0/(4A)$ lies below all collar and
mass-lemma thresholds. Suppose that a current interval $[l_j,u_j]$
contains $R$ and has width $w_j$. At its midpoint $a_j$ query with
\begin{equation}\label{eq:shrinking-layer-schedule}
            e_j=\frac{w_j}{4A},\qquad q=\frac34.
\end{equation}
For label one retain $[a_j-Ae_j,u_j]$; for label zero retain
$[l_j,a_j+Ae_j]$. Equation~\eqref{eq:shrinking-label-invariant}
proves that the retained interval still contains $R$. It is a subinterval
of the preceding interval and has width exactly $q w_j$, whichever
label occurs. Consequently
\[
                     w_j=w_0q^j.
\]
This deterministic width schedule holds also on an erroneous record;
only the containment assertion uses successful labels.

Put $h\asymp\nu^{1/(s-2)}$ and choose a positive dyadic final radial
tolerance $E\asymp c h^s$. Let $K$ be the least integer with
$w_K\le2E$.
For all sufficiently small $E$, one has $K\ge1$ and
\begin{equation}\label{eq:shrinking-final-width}
                   2qE<w_K\le2E,
             \qquad K\le C\log(C/E).
\end{equation}
The final midpoint has radial error at most $E$. Let $M_h\le C h^{-1}$
be a deterministic bound on the total number of radial searches,
including both scales when present. At level $j=0,\ldots,K-1$ of
each search use conditional failure allowance
\begin{equation}\label{eq:shrinking-confidence}
 \varepsilon_j=\frac{\delta}{4M_h}\,2^{-(K-1-j)}.
\end{equation}
The sum for one search is less than $\delta/(2M_h)$, so all fine
tests succeed, conditional on coarse success, with probability at least
$1-\delta/2$. This allocation permits relatively larger failure
allowances at the more expensive fine levels. Locations, labels and
the resulting subintervals may depend on the full past; fresh
preparations give the required conditional bounds at every level.

Write $\kappa=\gamma+3/2$. By Proposition~\ref{prop:rare-query},
the attempts in one search are at most
\begin{align*}
 C\sum_{j=0}^{K-1} e_j^{-\kappa}
                                  \log(C/\varepsilon_j)
 &\le C\sum_{j=0}^{K-1} w_j^{-\kappa}
          \left\{\log(CM_h/\delta)+(K-1-j)\log2\right\}\\
 &\le C w_K^{-\kappa}\log(CM_h/\delta).
\end{align*}
For the last inequality put $k=K-j$ and use the convergent sums
$\sum_{k\ge1}q^{\kappa k}$ and
$\sum_{k\ge1}kq^{\kappa k}$. Equation~\eqref{eq:shrinking-final-width}
therefore bounds the total fine cost by
\[
                         C M_h E^{-\kappa}\log(CM_h/\delta).
\]
The coarse cost is $C\log(C/\delta)$ and is absorbed by this bound.
Substitution of $M_h\le C h^{-1}$ and $E\asymp c h^s$ proves
\eqref{eq:shrinking-attempts}. There are at most two centers per
level and $M_hK$ levels in all; the coarse stencil count is fixed.
This proves \eqref{eq:shrinking-centers}.

All bracket endpoints, layer widths and updates are dyadic when $A$
is chosen as a power of two. For an explicit implementation, choose
one dyadic target mesh so that each nominal point is computed within
$E/(8A)$ of its ideal radial point. Since
$w_{K-1}>2E$, this error is less than $e_j/2$ at every level.
For sufficiently small $E$ it also lies within the fixed bracket
neighborhood of Lemma~\ref{lem:rare-coarse-normals}. Add the exact
dyadic compass shift after rounding the target once. Thus the preceding
radial conversion and normal certification apply to every physical
command. Certified conservative integer batch sizes preserve the
probability and cost bounds. The finest coordinates and all batch
descriptions require only
$O(\log(C/\nu)+\log(C/\delta))$ digits.

Finally, the radial errors are $O(h^s)$ on the common successful event.
The interpolation, support subtraction, convexity and period-locking
steps of Theorem~\ref{thm:rare-stationary} apply without change. At
two known scales, include all searches in $M_h$, use the same uniform
constants, and apply its complete-component correspondence and common
footprint reconstruction. This gives the stated output conclusions
with total failure probability at most $\delta$.
\end{proof}


\begin{corollary}[The stationary minimax power]
\label{cor:stationary-minimax}
Fix a bounded class $\mathcal B_0\subset\mathcal B_\eta$ with the
priors of Theorem~\ref{thm:rare-stationary} and containing the preceding
packing, a confidence
$0<\delta<1/4$, and the same known uniform-disk laws. Let
$N^*_{\rm pool}(\nu,\delta)$ be the infimum of the worst-case
expected number of attempts among controllers with uniform geometric
error at most $\nu$ and success probability at least $1-\delta$,
using the fixed pooled compass and one fixed radius. Define
$N^*_{\rm short}(\nu,\delta)$ in the same way, allowing arbitrary
directions, lengths at most $t_{\max}$, and the fixed radius interval.
Choose the pooled compass length below $t_{\max}$ and with
$2t+2R_+<d_0$. Then, with
$q_0=(3s/2+1)/(s-2)$,
\begin{equation}\label{eq:stationary-sharp-bracket}
 c\nu^{-q_0}
 \le N^*_{\rm short}(\nu,\delta)
 \le N^*_{\rm pool}(\nu,\delta)
 \le C\nu^{-q_0}\log\frac{C}{\nu\delta}.
\end{equation}
In particular both observation designs have the same stationary
minimax power at fixed confidence. The statement leaves the
logarithmic factor and its confidence dependence unresolved.
\end{corollary}
\begin{proof}
Theorem~\ref{thm:sharp-stationary-lower} gives the first inequality:
its centered-shape loss is weaker than the geometric loss here.
The second is inclusion of the command classes. The rare-collision
upper construction, with the decreasing-resolution allocation of
Proposition~\ref{prop:shrinking-rare}, gives the final inequality.
That construction has a deterministic attempt cap, which is also a
bound for the expected-cost criterion used in this corollary.
Thus the lower and upper assertions refer to the same accuracy,
confidence, physical class, and expected-cost criterion.
\end{proof}

The density-independent upper theorem applies to a larger nuisance
class than the fixed disk experiment used for this converse. The
minimax conclusion concerns $\gamma=0$ and fixed positive disk
spreads with the stated geometry. It does not assert the necessity
of the upper exponent for every $\gamma>0$, nor does it make its
constants uniform when a permitted spread tends to zero.
